% !TEX root = ../thesis.tex
% ==============================================================================
%  USER CUSTOMIZATION & CONFIGURATION LAYER (styles/custom.tex)
% ==============================================================================
%
%  WELCOME TO THE USER CUSTOMIZATION FILE!
%  This is the designated place for all your personal macros, custom packages,
%  math operators, visual environments, and layout overrides.
%
% ------------------------------------------------------------------------------
%  ARCHITECTURE & SEPARATION OF CONCERNS:
% ------------------------------------------------------------------------------
%  - styles/common.tex (Core Engine):
%      Houses the core packages, geometry, fonts, footnote engines, theorem
%      environments, and bidirectional layout machinery. DO NOT MODIFY common.tex
%      directly, so that template updates can be applied cleanly.
%
%  - styles/custom.tex (User Space - THIS FILE):
%      Houses your specific project settings, custom commands, theorem tweaks,
%      colors, and optional layout overrides. Edit this file freely.
%
% ------------------------------------------------------------------------------
%  GUIDE: WHAT YOU CAN CONFIGURE IN THIS FILE:
% ------------------------------------------------------------------------------
%  1. Math Operators & Notation:
%     Define custom operators using \newcommand or \DeclareMathOperator.
%     Always wrap LaTeX math in \ensuremath{...} or \operatorname{...} for safety.
%
%  2. Visual Elements & Highlight Boxes:
%     Create or modify tcolorbox environments (e.g., the sleek 'academicbox' below)
%     for callouts, case studies, or theorem highlights.
%
%  3. Title Page Spacing & Elastic Springs:
%     The frontmatter uses elastic vertical glue (springs with fill weights).
%     You can easily adjust the spacing above or below any element (author,
%     title, supervisor, logo scale) using the commented length overrides below.
%
%  4. Adding Extra Packages:
%     If you need extra packages (e.g., tikz, pgfplots, listings, minted):
%     Add them in Section 4 of this file. Note: Packages modifying bidirectional
%     text, fonts, or footnotes should be loaded with care to avoid conflicts
%     with XePersian.
% ==============================================================================


% ==============================================================================
% 1. PROJECT MATHEMATICAL OPERATORS & MACROS
% ==============================================================================
% Define your research-specific mathematical notations here.
% Using \ensuremath ensures the macro works cleanly in both text and math modes.

\newcommand{\OPT}{\ensuremath{\mathop{\mathrm{OPT}}}}
\newcommand{\APX}{\ensuremath{\mathop{\mathrm{APX}}}}
\newcommand{\ALG}{\ensuremath{\mathop{\mathrm{ALG}}}}

% Examples of additional operators you might add:
% \newcommand{\argmax}{\mathop{\mathrm{arg\,max}}}
% \newcommand{\argmin}{\mathop{\mathrm{arg\,min}}}
% \newcommand{\Expect}{\ensuremath{\mathbb{E}}}
% \newcommand{\Var}{\ensuremath{\mathrm{Var}}}
% \newcommand{\Loss}{\ensuremath{\mathcal{L}}}


% ==============================================================================
% 2. CUSTOM BOXES & VISUAL CALLOUT ENVIRONMENTS
% ==============================================================================
% Below is a formal, grayscale academic box (built with tcolorbox) ideal for
% summaries, problem statements, hypotheses, or key takeaways.

\newtcolorbox{academicboxenv}[2][]{%
    enhanced,
    breakable,
    colback=black!1.5!white,
    colframe=black!75,
    coltitle=white,
    fonttitle=\bfseries\footnotesize,
    title={#2},
    attach boxed title to top center={yshift*=-\tcboxedtitleheight/2},
    boxed title style={
        enhanced,
        colback=black!90,
        colframe=black!90,
        arc=2mm,
        boxrule=0pt,
        left=14pt,
        right=14pt,
        top=4pt,
        bottom=3pt
    },
    boxrule=0.8pt,
    borderline={0.3pt}{2pt}{black!20},
    left=16pt,
    right=16pt,
    top=11pt,
    bottom=13pt,
    before skip=18pt,
    after skip=18pt,
    #1
}

% Usage: \academicbox{Box Title}{Box Content}
\newcommand{\academicbox}[3][]{%
\begin{academicboxenv}[#1]{#2}
#3
\end{academicboxenv}%
}


% ==============================================================================
% 3. OPTIONAL SPACING & LAYOUT OVERRIDES
% ==============================================================================
% Default layout dimensions are established in styles/common.tex.
% If your university department requires slightly different margins or spacings,
% uncomment and adjust the relevant lines below:
%
% ------------------------------------------------------------------------------
% A. Persian & English Title Page Elastic Spacing (Spring Weights)
% ------------------------------------------------------------------------------
% The title page uses TeX 'fill' glues that automatically expand to balance
% the vertical content on an A4 page. Adjust the weights to shift elements up/down:
%
% \setlength{\thesistitlebeforedept}{0.1cm}                  % Space before department (FA)
% \setlength{\thesistitleenbeforedept}{0.2cm}                % Space before department (EN)
% \setlength{\thesistitlebeforedatatype}{0pt plus 1.5fill}    % Spring weight before degree/type
% \setlength{\thesistitlebeforetitle}{0pt plus 2.0fill}       % Spring weight before title
% \setlength{\thesistitlebeforeauthor}{0pt plus 2.2fill}      % Spring weight before author
% \setlength{\thesistitlebeforesupervisor}{0pt plus 1.0fill}  % Spring weight before supervisor
% \setlength{\thesistitleenbeforesupervisor}{0pt plus 1.0fill}% Spring weight before EN supervisor
% \setlength{\thesistitlebeforeadvisor}{0pt plus 0.8fill}     % Spring weight before advisor
% \setlength{\thesistitleenbeforeadvisor}{0pt plus 0.8fill}   % Spring weight before EN advisor
% \setlength{\thesistitlebeforeexaminer}{0pt plus 0.8fill}    % Spring weight before examiner
% \setlength{\thesistitleenbeforeexaminer}{0pt plus 0.8fill}  % Spring weight before EN examiner
% \setlength{\thesistitlebeforedate}{0pt plus 1.8fill}        % Spring weight before date
% \setlength{\thesistitleenbeforedate}{0pt plus 1.8fill}      % Spring weight before EN date
%
% --- University Logo Scale ---
% Adjust the size of the university logo on both Persian and English title pages:
% \renewcommand{\ThesisLogoScale}{0.25}                      % Default scale factor is 0.25
%
% ------------------------------------------------------------------------------
% B. Abstract & Frontmatter Typography
% ------------------------------------------------------------------------------
% Adjust vertical gaps or line-spread for abstract and acknowledgment pages:
% \setlength{\thesisabstractbeforespace}{0cm}                % Gap below Persian abstract title
% \setlength{\thesisabstractenbeforespace}{0.25cm}           % Gap below English abstract title
% \renewcommand{\thesisabstractenstretch}{0.73}              % English abstract line height ratio
% \setlength{\thesisackbeforespace}{1cm}                     % Gap below Acknowledgments title
%
% ------------------------------------------------------------------------------
% C. Footnote Spacing
% ------------------------------------------------------------------------------
% Gap between the bottom of the main text and the two-column footnotes:
% \setlength{\skip\footins}{1cm}                             % Default is 1cm

% ------------------------------------------------------------------------------
% D. Global Table Rules & Line Thickness
% ------------------------------------------------------------------------------
% Standardize table rule thickness across all tables to 0.5pt:
\setlength{\heavyrulewidth}{0.5pt}   % Default width for \toprule and \bottomrule
\setlength{\lightrulewidth}{0.5pt}   % Default width for \midrule
\setlength{\arrayrulewidth}{0.5pt}   % Default width for \vline, \hline, \hhline, and |


% ==============================================================================
% 4. ADDITIONAL USER PACKAGES (OPTIONAL)
% ==============================================================================
% If your research requires specialized packages (e.g., TikZ, listings, pgfplots),
% load them here:
%
% \usepackage{tikz}
% \usetikzlibrary{shapes,arrows,positioning}
% \usepackage{listings}
% \usepackage{pgfplots}
% \pgfplotsset{compat=1.18}
